
FALTARIA HACER CLARA LA DISTINCION ENTRE LOS DOS CAPITULOS EN SU SITIO CORRESPONDIENTE
\chapter{Results and combination with the Run-2 data set}
\label{ch:bbtt_results}



\section{Statistical procedure}
\label{sec:bbtt_statistics}


\begin{table}[tbp]
\centering
\small
\begin{tabular}{lllc}
\toprule
Source & Affected processes & Type & Correlation across periods \\
\midrule
Integrated luminosity                     & simulated              & normalization & partial \\
Signal cross section                      & signal                 & normalization & full    \\
Background cross sections                 & simulated              & normalization & full    \\
Branching fractions                       & with a Higgs boson     & normalization & full    \\
$b\bar{b}H$ normalization                 & $b\bar{b}H$            & normalization & full    \\
\midrule
Pile-up reweighting                       & simulated              & shape & full    \\
Trigger efficiency                        & simulated              & shape & none    \\
Muon and electron identification and isolation & simulated         & shape & full    \\
Electron energy scale and resolution      & simulated              & shape & full    \\
$\tau_h$ identification                   & simulated              & shape & mixed   \\
$\tau_h$ energy scale                     & simulated              & shape & none    \\
Jet energy scale                          & simulated              & shape & mixed   \\
Jet energy resolution                     & simulated              & shape & none    \\
$b$ tagging                               & simulated              & shape & mixed   \\
Acceptance (QCD scale and PDF)            & simulated              & shape & full    \\
\midrule
Drell--Yan correction, fit                & Drell--Yan             & shape & full    \\
Drell--Yan correction, statistical        & Drell--Yan             & shape & none    \\
Top quark $p_T$ reweighting               & $t\bar{t}$             & shape & full    \\
\midrule
Fake factor envelopes                     & \textit{fakes}         & shape & full    \\
Fake factor statistics                    & \textit{fakes}         & shape & none    \\
Template region statistics                & \textit{fakes}         & shape & none    \\
\midrule
Size of the simulated samples             & simulated              & shape & none    \\
\bottomrule
\end{tabular}
\caption[Summary of the systematic uncertainties]{%
    Summary of the sources of systematic uncertainty considered in the analysis, with the
    processes they affect, whether they modify only the normalization or also the shape of
    the distributions, and their correlation across data-taking periods.}
\label{tab:systematics_summary}
\end{table}


% Placed before the results, as in Jaime's thesis, so the results section can be
% read without interruption.

\subsection{The likelihood model}
% Binned maximum-likelihood fit over the 10 SRs x 5 data-taking periods,
% performed simultaneously. COMBINE / RooFit / RooStats.
% -> code: LimitsArea/, scripts/NBAtoInferenceDatacardsParser.py

\subsection{Observed and expected limits}
% Profile likelihood ratio test statistic; asymptotic modified frequentist CLs.
% Parameters of interest: mu, sigma(HH), or the coupling modifiers.

\subsection{Nuisance parameters}
% Normalisation vs shape uncertainties; correlation scheme across periods
% and across processes.

\section{Results}
\label{sec:bbtt_results}

\subsection{Post-fit distributions}
% DNN scores in the 10 SRs after the fit; show 2024 as the representative period
% (largest L_int), noting the others behave similarly.

\subsection{Signal strength and significance}
% mu = 2.6 +2.1 -1.7 ; observed (expected) significance 1.5 (0.7) sigma.
% No significant excess over the background-only hypothesis.

\subsection{Upper limits on the HH production cross section}
% Per category and combined:
%   inclusive ggHH+qqHH : 6.6 (3.9) x SM   -- resolved 2b drives it
%   qqHH only           : 81 (85)  x SM    -- VBF category drives it
% Per-year breakdown (Table 3): 2022 -> 17 (11), 2023 -> 8 (11), 2024 -> 8.3 (4.4).
% Explain why 2024 dominates: luminosity, new triggers, UParT.

\subsection{Constraints on the Higgs boson couplings}
%   -4.1 < kappa_lambda < 11.2  (expected -2.2 < kappa_lambda < 9.2)
%   -0.2 < kappa_2V < 2.3       (expected -0.2 < kappa_2V < 2.4)

\section{Combination with the Run-2 data set}
\label{sec:bbtt_combination}
% (3) RUN 2 ENTERS HERE, at datacard level.

\subsection{Combination procedure}
% Combined with the published Run-2 search in the same final state
% \cite{CMS-HIG-20-010}: 138/fb at 13 TeV -> 310/fb in total.
% - Some Run-2 theory uncertainties were updated after publication (an extra
%   uncertainty on heavy-flavour associated single-H production), so the standalone
%   Run-2 result differs slightly from the published one -- say this explicitly.
% - Correlation scheme: theory uncertainties (branching ratios, PDF, alpha_s,
%   scales, and notably the joint QCD+m_top uncertainty in ggHH) correlated across
%   Runs; all experimental uncertainties uncorrelated, because detector conditions
%   and calibrations differ.
% - The likelihood is the union of all binned shapes from both analyses; the
%   coupling parametrisation follows the same matrix-inversion method in both.

\subsection{Combined results}
% mu = 0.9 +1.6 -1.4 ; observed (expected) significance 0.7 (0.8) sigma.
%   inclusive : 4.0 (3.0) x SM   [Run 2 alone 3.5 (5.3) ; Run 3 alone 6.6 (3.9)]
%   qqHH only : 62 (70)  x SM
%   -2.5 < kappa_lambda < 9.4  (expected -1.5 < kappa_lambda < 8.4)
%   0.02 < kappa_2V < 2.1      (expected -0.1 < kappa_2V < 2.2)
% 2D exclusion in the kappa_2V - kappa_lambda plane, best fit at (7.1, 0.6).
% Worth a comment: Run 3 observed is weaker than expected while Run 2 is stronger
% than expected -- explain how the two combine rather than leaving it unremarked.

\section{Comparison with other HH searches}
\label{sec:bbtt_comparison}
% Borrowed from Jaime (his Sec. 5.5), and it earns its place.
% CMS Run-2 combination 3.4 (2.5); ATLAS 2.9 (2.4); ATLAS+CMS 2.5 (1.7 / 2.8);
% ATLAS bbtautau with 196/fb -> 5.3 (2.4 / 2.8). Where bbtautau sits among
% bbgammagamma, bbbb and bbWW.

\begin{table}[t]
\centering
\begin{tabular}{llccccl}
\toprule
 & Dataset & $\mathcal{L}$ [fb$^{-1}$] & Obs. & Exp. & Exp.$^{\dagger}$ & Source \\
\midrule
\multirow{3}{*}{CMS}
  & Run 2 (2016--2018)         & 138 & 3.5 & 5.3 & 5.30 & \cite{CMS-PAS-HIG-25-008} \\
  & Run 3 (2022--2024)         & 172 & 6.6 & 3.9 & 4.35 & \cite{CMS-PAS-HIG-25-008} \\
  & Run 2 + Run 3              & 310 & 4.0 & 3.0 & 4.50 & \cite{CMS-PAS-HIG-25-008} \\
\midrule
\multirow{4}{*}{ATLAS}
  & Run 2, original            & 140 & 5.9 & 3.3 & 3.32 & \cite{ATLAS-HH-bbtautau-Run2} \\
  & Run 2, reanalysed          & 140 & 5.4 & 2.6 & 2.62 & \cite{ATLAS-HH-bbtautau-Run3} \\
  & Run 3 (2022--2023)         &  56 & 7.9 & 4.3 & 2.74 & \cite{ATLAS-HH-bbtautau-Run3} \\
  & Run 2 + Run 3              & 196 & 5.3 & 2.4 & 2.86 & \cite{ATLAS-HH-bbtautau-Run3} \\
\bottomrule
\end{tabular}
\caption[Comparison of CMS and ATLAS upper limits in HH $\to$ bb$\tau\tau$]{%
Observed and expected 95\% CL upper limits on the inclusive HH production cross
section relative to the SM prediction, in the bb$\tau\tau$ final state. Expected
limits are quoted under the background-only ($\mu = 0$) hypothesis. The CMS
entries are the inclusive ggHH\,$+$\,qqHH limits of Table~3 of
Ref.~\cite{CMS-PAS-HIG-25-008}; the ATLAS Run~3 entries are those of Table~4 of
Ref.~\cite{ATLAS-HH-bbtautau-Run3}, where ``reanalysed'' denotes the same Run~2
dataset processed with the improved Run~3 analysis, to be compared with the
originally published result of Ref.~\cite{ATLAS-HH-bbtautau-Run2}.
$^{\dagger}$Expected limit rescaled to a common integrated luminosity of
138~fb$^{-1}$ as $\mu\sqrt{\mathcal{L}/138~\mathrm{fb}^{-1}}$, i.e.\ assuming
purely statistical scaling; this column is derived here and is not quoted in the
references. Note that the two Run~3 datasets are not equivalent, as the ATLAS
result does not yet include the 2024 data-taking period.}
\label{tab:hh_bbtautau_limits}
\end{table}

\section{Summary and outlook}
\label{sec:bbtt_summary}
% Each of the reference theses closes its analysis chapter with a summary.
% - What was measured and with what sensitivity.
% - What limits it now: statistically dominated (74%), fake factors the leading
%   experimental systematic (25%).
% - The impact of the Chapter \ref{ch:taggers} algorithms on this result -- this is
%   where the thesis closes its own loop.
% - Outlook: remaining Run-3 luminosity, then HL-LHC.
SE ME HACE UN CAPITULO MUY LARGO (HAY 14 SECCIONES...) SE PUEDE SEPARAR EN VARIOS CAPITULOS? HACIENDOLO UNA "PARTE". ALGO COMO LO QUE TIENE JAIME? TENDRÍA SENTIDO ESA SEPARACION?


\chapter{Results and combination with the Run-2 data set}
\label{ch:bbtt_results_chapter}
% Entradilla: recuperar el hilo del capítulo anterior en dos o tres frases.
% Construida la predicción, comprobada frente a los datos y definido el
% discriminante, queda someterlo al contraste estadístico y extraer la medida.


\section{The statistical model of the search}
\label{sec:bbtt_statistics}
% SOLO lo propio de este análisis. El marco general -- verosimilitud, parámetros
% de estorbo y sus términos de restricción, cociente de verosimilitudes perfilado,
% CLs -- está en la Sección \ref{sec:stat_framework} del Capítulo III.
%
% 1) El ajuste: binado, de máxima verosimilitud, simultáneo sobre las diez
%    categorías x cinco periodos = cincuenta regiones de señal (Sección \ref{sec:bbtt_flow}).
%    Observable en cada una: la salida de la DNN de la Sección \ref{sec:bbtt_extraction},
%    con el binado adaptativo allí descrito.
% 2) Parámetro de interés: la intensidad de señal. Para las interpretaciones en
%    kappa_lambda y kappa_2V, el modelo que parametriza la dependencia de la
%    sección eficaz y de la cinemática (mismo método de inversión de matriz de
%    la Sección \ref{sec:bbtt_samples}).
% 3) Los parámetros de estorbo y su correlación: aquí va la Tabla
%    \ref{tab:systematics_summary}, que ahora está en el Capítulo VI.
%    Correlación entre periodos, y descorrelación por categoría y canal de
%    CMS_FF_stat y CMS_FF_template_stat.
% 4) La estadística de las muestras simuladas: el párrafo de autoMCStats
%    (umbral de diez sucesos efectivos, un parámetro por intervalo compartido
%    por encima y uno por proceso por debajo, n_eff = (sum w)^2 / sum w^2).
% -> code: LimitsArea/, scripts/NBAtoInferenceDatacardsParser.py,
%    config/run3_base_config.py (alias, syst_runPeriod_tag),
%    LimitsArea/scripts/SymmetriseAndSmoothShapes.py (renombrado de nuisances)
%
% MOVER AQUÍ la Tabla tab:systematics_summary desde el Capítulo VI.


\section{Results}
\label{sec:bbtt_results}

\subsection{Post-fit distributions}

\subsection{Signal strength and significance}

\subsection{Impact of the systematic uncertainties}
\label{subsec:bbtt_impacts}
% Prometido desde la entradilla de las sistemáticas (Cap. VI, l. 984).
% - Desglose por familia: Tabla 2 del PAS. Estadística total 74%; fake factors 25%;
%   etiquetado de b 10%; leptones y disparo 9%; energía de jets 7%; otras 3%;
%   teóricas HH 31%, single-H 27%, fondos 14%; variaciones intervalo a intervalo 35%.
%   Decir que se combinan en cuadratura y que por eso no suman 100%.
% - Pull and impact: ordena parámetros individuales por Delta(mu); complementa
%   la tabla, que ordena familias enteras.
% -> code: LimitsArea/scripts/run_syst_breakdown.sh (grupos theoHH, theoH,
%    theoAcc, theoB y las rgx por objeto)

\subsection{Upper limits on the HH production cross section}

\subsection{Constraints on the Higgs boson couplings}


\section{Combination with the Run-2 data set}
\label{sec:bbtt_combination}

\subsection{Combination procedure}

\subsection{Combined results}


\section{Comparison with other HH searches}
\label{sec:bbtt_comparison}


\section{Summary and outlook}
\label{sec:bbtt_summary}
